\documentclass[11pt]{article}

\usepackage[a4paper,margin=28mm]{geometry}
\usepackage[T1]{fontenc}
\usepackage[utf8]{inputenc}
\usepackage{lmodern}
\usepackage{microtype}
\usepackage{amsmath,amssymb,amsthm,mathtools}
\usepackage{graphicx}
\usepackage{enumitem}
\usepackage[hidelinks]{hyperref}
\usepackage{bookmark}

\setlist[itemize]{leftmargin=2em,itemsep=0.25em,topsep=0.4em}
\setlist[enumerate]{leftmargin=2.2em,itemsep=0.25em,topsep=0.4em}
\allowdisplaybreaks

\newtheorem{theorem}{Theorem}[section]
\newtheorem{proposition}[theorem]{Proposition}
\newtheorem{lemma}[theorem]{Lemma}
\newtheorem{corollary}[theorem]{Corollary}
\newtheorem{question}[theorem]{Question}
\theoremstyle{definition}
\newtheorem{definition}[theorem]{Definition}
\newtheorem{remark}[theorem]{Remark}

\newcommand{\N}{\mathbb N}
\newcommand{\R}{\mathbb R}
\newcommand{\dd}{d_2}
\newcommand{\BM}{d_{\mathrm{BM}}}
\newcommand{\Id}{\operatorname{Id}}

\title{Question 6.17(2) on Unconditional FDDs\\
\large Audit, literature status, and an exact tail-localization criterion}
\author{Research report prepared from the supplied paper and partial solution}
\date{22 June 2026}

\begin{document}
\maketitle

\begin{abstract}
The question is whether every non-Hilbertian Banach space with an unconditional
finite-dimensional decomposition (UFDD) contains a non-Hilbertian subspace with
an unconditional basis.  I did not locate a published proof or counterexample.
The proposed sufficient condition is correct after repairing one uniformity
argument, but its description in terms of Pisier's property (H) is inaccurate.
More significantly, I prove a converse: for a fixed UFDD, the proposed
finite-dimensional tail-witness condition is equivalent to the desired
conclusion.  Thus the partial result is an exact localization of the open
problem rather than a strictly weaker theorem.  I also reduce the hard case to
a non-Hilbertian Hilbertian sum $(\bigoplus E_n)_2$ of finite-dimensional
spaces and prove several positive cases.  In particular, any counterexample
must have type $p$ for every $p<2$ and cotype $q$ for every $q>2$.
\end{abstract}

\section{The source question and the verdict}

The source is W.~Cuellar Carrera, N.~de Rancourt and V.~Ferenczi,
\emph{Local Banach-space dichotomies and ergodic spaces}~\cite{CDF}.
Question~6.17(2) asks:
\begin{quote}
Does every non-Hilbertian space with an unconditional FDD have a
non-Hilbertian subspace with an unconditional basis?
\end{quote}

\begin{center}
\includegraphics[width=\textwidth]{figures/open_problem_crop.png}
\end{center}

The bibliographic attribution in the supplied packet is incorrect: the paper
is not by Ferenczi--Lopez-Abad--Mbombo--Todorcevic.  It appeared in the
\emph{Journal of the European Mathematical Society} in 2023; the preprint is
arXiv:2005.06458.

The outcome of the investigation is as follows.
\begin{enumerate}
\item The theorem in the packet is correct, but its proof of a uniform
      $\ell_2$ estimate has a logical gap.  A corrected tail version is proved
      below and suffices for the concatenation argument.
\item The phrase ``uniform tail property-(H) failure'' should be avoided.  The
      condition concerns Banach--Mazur distance of finite-dimensional spaces,
      while Pisier's property (H) controls equal-coefficient sums of
      unconditional sequences.  They are not equivalent.  Non-Hilbertian weak
      Hilbert spaces with unconditional bases show that the proposed tail
      condition can coexist with property (H)~\cite{Pisier,NTJ,ABR}.
\item Exact-wording searches, citation trails, and the classical and recent
      literature on UFDDs, weak Hilbert spaces, twisted Hilbert spaces and
      $d_2$-HI spaces did not reveal a full answer or counterexample through
      the search date.
\item Relative to a fixed UFDD, the packet's hypothesis is not merely
      sufficient: it is equivalent to the existence of the desired subspace.
      This exact tail-localization is the main additional result below.
\end{enumerate}

Throughout, spaces are real Banach spaces.  For a finite-dimensional space
$F$, write
\[
  \dd(F)=\BM(F,\ell_2^{\dim F}).
\]
A normalized basic sequence is $K$-unconditional when arbitrary sign changes
have norm at most $K$ times the original norm.

\section{A finite-dimensional estimate}

\begin{lemma}[Unconditional Euclidean comparison]\label{lem:finite-euclidean}
Let $G$ be $m$-dimensional and let $(x_i)_{i=1}^m$ be a normalized
$K$-unconditional basis of $G$.  For every $D>\dd(G)$, the sequence $(x_i)$ is
$KD$-equivalent to the unit vector basis of $\ell_2^m$:
\[
  \frac{1}{KD}\Big(\sum_{i=1}^m |a_i|^2\Big)^{1/2}
  \leq \Big\|\sum_{i=1}^m a_i x_i\Big\|
  \leq KD\Big(\sum_{i=1}^m |a_i|^2\Big)^{1/2}.
\]
\end{lemma}

\begin{proof}
Choose $T:G\to\ell_2^m$ with $\|T\|\,\|T^{-1}\|<D$.  Put
$x=\sum_i a_i x_i$ and let $(\varepsilon_i)$ be independent signs.
Unconditionality gives
\[
 K^{-1}\|x\|\leq
 \Big(\mathbb E\big\|\sum_i\varepsilon_i a_i x_i\big\|^2\Big)^{1/2}
 \leq K\|x\|.
\]
Orthogonality in $\ell_2^m$ gives
\[
 \frac{1}{\|T\|\|T^{-1}\|}\|a\|_2
 \leq
 \Big(\mathbb E\big\|\sum_i\varepsilon_i a_i x_i\big\|^2\Big)^{1/2}
 \leq \|T\|\|T^{-1}\|\|a\|_2.
\]
Combining the inequalities proves the assertion.
\end{proof}

\begin{remark}\label{rem:H-name}
Large $\dd(G)$ need not force the equal-sum vector of a particular
unconditional basis of $G$ to have abnormal size.  Hence the packet's
condition should be renamed rather than identified with failure of property
(H).
\end{remark}

\section{The corrected sufficient theorem}

Let $(E_n)$ be a UFDD of $X$, and write $P_{[r,s]}$ for its interval
projections.

\begin{definition}[Tail-witness property]
The UFDD $(E_n)$ has the \emph{tail-witness property} if there is
$\lambda\geq1$ such that, for every $N\in\N$ and $M\geq1$, there are
$N\leq r\leq s$ and
\[
  G\subseteq E_r\oplus\cdots\oplus E_s
\]
with a normalized $\lambda$-unconditional basis and $\dd(G)>M$.
\end{definition}

The next lemma repairs the gap in the packet.

\begin{lemma}[Uniformity on a tail]\label{lem:tail-uniformity}
Assume every normalized block sequence of $(E_n)$ spans a Hilbertian space.
Then there are $N_0\in\N$ and $C\geq1$ such that every normalized block
sequence supported after $N_0$ is $C$-equivalent to the unit vector basis of
$\ell_2$.
\end{lemma}

\begin{proof}
Suppose not.  Recursively choose finite normalized block sequences
\[
  (u_i^{(j)})_{i=1}^{m_j},\qquad j=1,2,\ldots,
\]
with successive, mutually disjoint supports, such that the $j$th sequence is
not $j$-equivalent to the unit vector basis of $\ell_2^{m_j}$.  Their
concatenation is a normalized block sequence.  It is unconditional and, by
hypothesis, spans a Hilbertian space.  A normalized unconditional basic
sequence in a Hilbert space is equivalent to the $\ell_2$ basis with one fixed
constant.  The same constant works on every finite consecutive piece, which
contradicts the choice of the pieces.
\end{proof}

The packet asserted a global constant and argued that failure of global
uniformity supplies bad sequences in every tail.  That implication was not
proved.  Lemma~\ref{lem:tail-uniformity} is the direct gliding-hump conclusion
and is all that the proof needs.  A global estimate can subsequently be
recovered by separating the uniformly bounded finite head from the controlled
tail, but that extra step was absent from the packet.

\begin{theorem}[Corrected packet theorem]\label{thm:sufficient}
Let $X$ have a UFDD $(E_n)$.  If $(E_n)$ has the tail-witness property, then
$X$ contains a non-Hilbertian subspace with an unconditional basis.
\end{theorem}

\begin{proof}
If some normalized block sequence spans a non-Hilbertian space, it already is
an unconditional basic sequence.  Otherwise take $N_0,C$ from
Lemma~\ref{lem:tail-uniformity}.

Choose successive interval-supported spaces $G_k$, all after $N_0$, such that
$\dd(G_k)\to\infty$ and each $G_k$ has a normalized
$\lambda$-unconditional basis $(x_i)_{i_k\leq i<i_{k+1}}$.  For finitely
supported scalars $(a_i)$ and signs $(\varepsilon_i)$, put
\[
 u_k=\sum_{i_k\leq i<i_{k+1}}a_i x_i,
 \qquad
 v_k=\sum_{i_k\leq i<i_{k+1}}\varepsilon_i a_i x_i,
\]
and let $b_k=\|u_k\|$, $c_k=\|v_k\|$.  When nonzero, write
$y_k=u_k/b_k$ and $z_k=v_k/c_k$.  Internal unconditionality gives
$c_k\leq\lambda b_k$, while $(y_k)$ and $(z_k)$ are normalized FDD block
sequences after $N_0$.  Hence
\begin{align*}
 \Big\|\sum_i\varepsilon_i a_i x_i\Big\|
 &=\Big\|\sum_k c_k z_k\Big\| \\
 &\leq C\Big(\sum_k c_k^2\Big)^{1/2}
 \leq \lambda C\Big(\sum_k b_k^2\Big)^{1/2} \\
 &\leq \lambda C^2\Big\|\sum_k b_k y_k\Big\|
 =\lambda C^2\Big\|\sum_i a_i x_i\Big\|.
\end{align*}
Thus the concatenation is unconditional.  Its closed span contains every
$G_k$ and therefore cannot be Hilbertian, since $\dd(G_k)\to\infty$.
\end{proof}

\section{A converse: exact tail localization}

The converse uses a finite-rank avoidance lemma.

\begin{lemma}[Annihilating a finite-rank operator]\label{lem:finite-rank}
Let $Y$ be a non-Hilbertian space with a normalized unconditional basis.
There is $L\geq1$, depending only on that basis, such that for every
finite-rank $T:Y\to F$, every $\eta>0$, and every $M\geq1$, there is a
finite-dimensional $G\subseteq Y$ satisfying:
\begin{enumerate}[label=(\roman*)]
\item $G$ has a normalized $L$-unconditional basis;
\item $\dd(G)>M$;
\item $\|T|_G\|<\eta$.
\end{enumerate}
\end{lemma}

\begin{proof}
Let $(y_j)$ be a normalized $K$-unconditional basis of $Y$.  We use the
standard fact that a nonreflexive space with an unconditional basis contains a
normalized block sequence equivalent to the canonical basis of either $c_0$ or
$\ell_1$~\cite{AK}.  Every normalized block sequence of $(y_j)$ is
$K$-unconditional, so $L=K$ works in all cases.

\emph{Reflexive case.}
The basis is shrinking.  For the tail projection
$Q_m:Y\to[y_j:j\geq m]$, one has $\|TQ_m\|\to0$ because $T$ has finite rank.
Every basis tail is non-Hilbertian.  Consequently
\[
 \dd([y_j:m\leq j\leq r]),\qquad r\geq m,
\]
is unbounded: otherwise Lemma~\ref{lem:finite-euclidean} would make the
entire tail basis uniformly equivalent to $\ell_2$.  Choose $m$ with
$\|TQ_m\|<\eta$, then choose a coordinate interval $G$ in that tail with
$\dd(G)>M$.

\emph{$c_0$ case.}
Fix a normalized block sequence $(w_j)$ which is $D$-equivalent to the
$c_0$ basis.  It is weakly null, so $\|Tw_j\|\to0$.  Pass to a subsequence
with
\[
  \sum_j\|Tw_j\|<\eta/D.
\]
For $x=\sum_{j=1}^m a_jw_j$ we have $\max_j|a_j|\leq D\|x\|$, and hence
$\|Tx\|<\eta\|x\|$.  For $G_m=[w_1,\ldots,w_m]$,
\[
  \dd(G_m)\geq D^{-2}\BM(\ell_\infty^m,\ell_2^m)
             =D^{-2}\sqrt m.
\]
Take $m$ large.

\emph{$\ell_1$ case.}
Fix a normalized block sequence $(v_j)$ which is $D$-equivalent to the
$\ell_1$ basis.  Since $F$ is finite-dimensional, $(Tv_j)$ has a Cauchy
subsequence.  Choose successive disjoint pairs from it so rapidly that, after
normalization,
\[
 w_i=\frac{v_{p_{2i}}-v_{p_{2i-1}}}
          {\|v_{p_{2i}}-v_{p_{2i-1}}\|},
 \qquad
 \sum_i\|Tw_i\|<\eta/D_1,
\]
where $(w_i)$ is $D_1$-equivalent to the $\ell_1$ basis for a fixed $D_1$.
Then $\|T|_{[w_1,\ldots,w_m]}\|<\eta$ and
\[
 \dd([w_1,\ldots,w_m])
 \geq D_1^{-2}\BM(\ell_1^m,\ell_2^m)
 =D_1^{-2}\sqrt m.
\]
Again take $m$ large.
\end{proof}

\begin{theorem}[Exact tail-localization criterion]\label{thm:equivalence}
Let $X$ have a UFDD $(E_n)$.  The following are equivalent.
\begin{enumerate}[label=(\alph*)]
\item $X$ contains a non-Hilbertian subspace with an unconditional basis.
\item $(E_n)$ has the tail-witness property.
\end{enumerate}
\end{theorem}

\begin{proof}
The implication (b)$\Rightarrow$(a) is Theorem~\ref{thm:sufficient}.

For the converse, let $Y\subseteq X$ be non-Hilbertian with an unconditional
basis.  Fix $N$ and apply Lemma~\ref{lem:finite-rank} to
$T=P_{[1,N-1]}|_Y$.  Let $\delta=1/8$.  For a number $R$ to be chosen, obtain
$G\subseteq Y$ with a normalized $L$-unconditional basis, $\dd(G)>R$, and
\[
 \|P_{[1,N-1]}|_G\|<\delta.
\]
Because $G$ is finite-dimensional, choose $s\geq N$ so that
$\|P_{(s,\infty)}|_G\|<\delta$.  Put $S=P_{[N,s]}|_G$.  Then
$\|S-\Id_G\|<2\delta<1$, so $S$ is an isomorphism onto its image and
\[
 \|S\|\,\|S^{-1}\|
 \leq \frac{1+2\delta}{1-2\delta}=:D_0.
\]
The image $SG$ lies in $E_N\oplus\cdots\oplus E_s$.  After normalizing the
image basis, it is $LD_0$-unconditional, and
\[
  \dd(SG)\geq \frac{\dd(G)}{D_0}.
\]
Choosing $R>D_0M$ proves the tail-witness property with the fixed constant
$\lambda=LD_0$.
\end{proof}

\begin{remark}
The packet's hypothesis is therefore exactly the desired conclusion localized
into arbitrarily remote finite intervals of the fixed UFDD.  This validates the
mechanism, but also shows that it does not by itself reduce the open problem to
a genuinely weaker statement.
\end{remark}

\section{The hard case: a Hilbertian sum of bad blocks}

\begin{theorem}[Block/Hilbertian-sum reduction]\label{thm:hardcase}
Let $X$ be a non-Hilbertian space with a UFDD $(E_n)$.  Then one of the
following holds.
\begin{enumerate}[label=(\roman*)]
\item Some normalized block sequence of $(E_n)$ spans a non-Hilbertian space;
      hence the question has a positive answer for $X$.
\item There are $N_0$ and $C$ such that, for all finitely supported
      $x_n\in E_n$ with $n\geq N_0$,
      \[
        C^{-1}\Big(\sum_{n\geq N_0}\|x_n\|^2\Big)^{1/2}
        \leq \Big\|\sum_{n\geq N_0}x_n\Big\|
        \leq
        C\Big(\sum_{n\geq N_0}\|x_n\|^2\Big)^{1/2}.
      \]
      Consequently the tail is isomorphic to
      $(\bigoplus_{n\geq N_0}E_n)_2$.  Moreover, for every $N\geq N_0$,
      \[
        \sup_{n\geq N}\dd(E_n)=\infty.
      \]
\end{enumerate}
\end{theorem}

\begin{proof}
If (i) fails, Lemma~\ref{lem:tail-uniformity} gives $N_0,C$.  Apply its
$C$-equivalence estimate to the normalized one-block vectors
$x_n/\|x_n\|$ to obtain the displayed inequality.

If $\sup_{n\geq N}\dd(E_n)<\infty$ on some tail, choose uniformly controlled
isomorphisms of those $E_n$ onto Euclidean spaces and take their Hilbertian
direct sum.  The displayed estimate would make that tail Hilbertian.  Adding
finitely many earlier summands would then make $X$ Hilbertian, a contradiction.
\end{proof}

In case (ii), the tail is reflexive and $\ell_2$-saturated.  The unresolved
issue is internal to the finite-dimensional blocks: their Euclidean distortion
is unbounded, but uniformly unconditional non-Euclidean subspaces might become
harder to find farther out.

For $\lambda\geq1$ and $N\in\N$, define
\[
 \Phi_\lambda(N)=\sup\left\{\dd(G):
 \begin{array}{l}
 G\subseteq E_r\oplus\cdots\oplus E_s,\quad N\leq r\leq s,\\
 G\text{ has a normalized }\lambda\text{-unconditional basis}
 \end{array}\right\}.
\]
By Theorem~\ref{thm:equivalence}, the answer for $X$ is positive exactly when
$\Phi_\lambda(N)=\infty$ for some fixed $\lambda$ and every $N$.

\section{Concrete positive consequences}

\begin{corollary}[Bounded-dimensional UFDD]\label{cor:bounded-dim}
If $X$ is non-Hilbertian and has a UFDD $(E_n)$ with
$\sup_n\dim E_n<\infty$, then $X$ contains a non-Hilbertian block subspace
with an unconditional basis.
\end{corollary}

\begin{proof}
John's theorem gives $\dd(E_n)\leq\sqrt{\dim E_n}$ uniformly~\cite{AK}.
Alternative (ii) of Theorem~\ref{thm:hardcase} is impossible.
\end{proof}

Thus classical examples with a two-dimensional UFDD but no unconditional basis
for the whole space, including the Kalton--Peck and Komorowski-type examples,
are not counterexamples to the subspace question~\cite{CK,KP,Komorowski}.

\begin{corollary}[Uniformly unconditional blocks]\label{cor:uniform-blocks}
Suppose $X$ is non-Hilbertian with a UFDD $(E_n)$, and every $E_n$ has a
normalized $K$-unconditional basis for one fixed $K$.  Then $X$ contains a
non-Hilbertian subspace with an unconditional basis.
\end{corollary}

\begin{proof}
In alternative (i) of Theorem~\ref{thm:hardcase} we are done.  In alternative
(ii), the numbers $\dd(E_n)$ are unbounded on every tail, so the individual
blocks furnish the tail-witness property with $\lambda=K$.
\end{proof}

\begin{corollary}[Uniformly Euclidean blocks]
Suppose $X$ has a UFDD $(E_n)$ satisfying
\[
  \sup_n\dd(E_n)<\infty.
\]
If $X$ is non-Hilbertian, then it contains a non-Hilbertian block subspace
with an unconditional basis.
\end{corollary}

\begin{proof}
Alternative (ii) of Theorem~\ref{thm:hardcase} is impossible.
\end{proof}

\begin{corollary}[Type--cotype obstruction]\label{cor:type-cotype}
For a Banach space $X$, put
\[
 p_X=\sup\{p\in[1,2]:X\text{ has type }p\},\qquad
 q_X=\inf\{q\in[2,\infty):X\text{ has cotype }q\},
\]
with $q_X=\infty$ if $X$ has no finite cotype.  Let $X$ be non-Hilbertian and
have a UFDD.  If $p_X<2$ or $q_X>2$, then $X$ contains a non-Hilbertian
subspace with an unconditional basis.  Consequently, every counterexample
would satisfy
\[
  p_X=q_X=2;
\]
equivalently, it would have type $p$ for every $p<2$ and cotype $q$ for every
$q>2$, without necessarily having type $2$ or cotype $2$.
\end{corollary}

\begin{proof}
Finite-dimensional perturbations do not change type or cotype exponents, so
every UFDD tail has the same indices.  The Maurey--Pisier theorem says that
$\ell_{p_X}$ is finitely representable in every tail and, if $q_X>2$,
$\ell_{q_X}$ is finitely representable there; when $q_X=\infty$, use
$\ell_\infty$~\cite{MP}.

Assume $p_X<2$.  For every tail and $m$, choose a subspace $H$ which is
$2$-isomorphic to $\ell_{p_X}^m$.  Since $H$ is finite-dimensional, a long
enough interval projection is a $2$-isomorphism on $H$.  Its image $G$ is
supported in one finite tail interval, has a normalized unconditional basis
with a universal constant, and
\[
 \dd(G)\geq \frac14\BM(\ell_{p_X}^m,\ell_2^m)
          =\frac14m^{1/p_X-1/2}\longrightarrow\infty.
\]
Thus the tail-witness property holds.  The case $2<q_X<\infty$ is identical,
using $\BM(\ell_{q_X}^m,\ell_2^m)=m^{1/2-1/q_X}$, and for $q_X=\infty$ use
$\BM(\ell_\infty^m,\ell_2^m)=\sqrt m$.
\end{proof}

Casazza--Kalton proved that a bounded-dimensional UFDD together with
Gordon--Lewis local unconditional structure yields an unconditional basis for
the whole space~\cite{CK}.  Corollary~\ref{cor:bounded-dim} assumes no local
unconditional structure, but concludes only for a suitable non-Hilbertian
subspace.

\section{Literature status}

The source paper proves a positive answer for $d_2$-minimal spaces
(Theorem~6.15 of~\cite{CDF}).  Its proof obtains uniformly unconditional,
increasingly non-Euclidean finite-dimensional witnesses from $d_2$-minimality
and failure of property (H), then uses essentially the concatenation mechanism
of the supplied packet.

Casazza--Kalton~\cite{CK} is the closest classical paper on assembling bases
inside UFDD summands.  It does not settle the present subspace problem for
unbounded block dimensions.  Kalton proved that a twisted Hilbert space with
an unconditional basis is Hilbertian, and his argument also controls
complemented unconditional basic sequences~\cite{Kalton}; this does not rule
out arbitrary uncomplemented subspaces of a twisted Hilbert space.

The 2025 paper of Argyros--Manoussakis--Motakis constructs classical models
related to the announced $\ell_2$-saturated $d_2$-HI problem, but explicitly
leaves the exotic construction for future work~\cite{AMM}.  The recent work of
de Rancourt--Kurka proves ergodicity for non-Hilbertian Orlicz sequence spaces
and several twisted Hilbert spaces, not Question~6.17(2)~\cite{RK}.

Exact searches for the wording of Question~6.17(2), searches combining
``non-Hilbertian'', ``unconditional FDD'', ``unconditional basis'',
``property (H)'', and ``twisted Hilbert'', and inspection of the citation trail
of the JEMS paper did not reveal a full solution.  The responsible conclusion
is therefore that no full answer was located and the question appears to
remain open.  This is a literature-search conclusion, not a claim about
unpublished work.

\section{What a counterexample must accomplish}

By Theorems~\ref{thm:equivalence} and~\ref{thm:hardcase}, after deleting
finitely many coordinates a counterexample may be assumed to be
\[
  X=\Big(\bigoplus_{n=1}^\infty E_n\Big)_2,
\]
where $\sup_n\dd(E_n)=\infty$.  By Corollary~\ref{cor:type-cotype}, it must
also satisfy $p_X=q_X=2$.  Nevertheless, for every fixed $\lambda$, some tail
must impose a uniform upper bound on $\dd(G)$ for every finite-dimensional
interval-supported $G$ with a normalized $\lambda$-unconditional basis.

Thus choosing individual blocks with large unconditional-basis constants is
not enough: one must prevent uniformly unconditional, non-Euclidean subspaces
from emerging across several successive blocks of the Hilbertian sum.  Random
finite-dimensional spaces and spaces with poor local unconditional structure
are natural candidates, but stability of the required obstruction under
finite Hilbertian sums is the missing step.

Conversely, a positive solution is equivalent to the following concrete
finite-dimensional principle.

\begin{question}[Residual finite-dimensional principle]\label{q:residual}
Let $(E_n)$ be finite-dimensional spaces such that
$(\bigoplus_n E_n)_2$ is non-Hilbertian.  Must there exist a constant
$\lambda$ and successive finite intervals $I_k$ containing subspaces
\[
  G_k\subseteq\Big(\bigoplus_{n\in I_k}E_n\Big)_2
\]
with normalized $\lambda$-unconditional bases and $\dd(G_k)\to\infty$?
\end{question}

Question~\ref{q:residual} is equivalent to Question~6.17(2) through the block
reduction and the exact tail-localization theorem.  It is the sharp endpoint
reached here.

\section{Conclusion}

The automatic solver found a valid concatenation mechanism.  Its theorem is
sound after two revisions: replace the flawed global-uniformity argument by
Lemma~\ref{lem:tail-uniformity}, and replace the property-(H) terminology by
the tail-witness terminology.  The principal extension is
Theorem~\ref{thm:equivalence}, which proves that the packet's condition is
necessary as well as sufficient.  The hard case is reduced to a
finite-dimensional extraction problem in non-Hilbertian Hilbertian sums, and
several substantial classes are ruled out as counterexamples.  I did not
obtain a complete proof or a counterexample.

\begin{thebibliography}{99}

\bibitem{AK}
F.~Albiac and N.~J. Kalton,
\emph{Topics in Banach Space Theory}, 2nd ed., Graduate Texts in Mathematics
233, Springer, 2016.

\bibitem{ABR}
S.~A. Argyros, K.~Beanland and T.~Raikoftsalis,
A weak Hilbert space with few symmetries,
\emph{C. R. Math. Acad. Sci. Paris} \textbf{348} (2010), 1293--1296.

\bibitem{AMM}
S.~A. Argyros, A.~Manoussakis and P.~Motakis,
Variants of the James tree space,
\emph{Israel J. Math.} \textbf{265} (2025), 1--77.

\bibitem{CK}
P.~G. Casazza and N.~J. Kalton,
Unconditional bases and unconditional finite-dimensional decompositions in
Banach spaces,
\emph{Israel J. Math.} \textbf{95} (1996), 349--373.

\bibitem{CDF}
W.~Cuellar Carrera, N.~de Rancourt and V.~Ferenczi,
Local Banach-space dichotomies and ergodic spaces,
\emph{J. Eur. Math. Soc.} \textbf{25} (2023), 3537--3598.

\bibitem{KP}
N.~J. Kalton and N.~T. Peck,
Twisted sums of sequence spaces and the three-space problem,
\emph{Trans. Amer. Math. Soc.} \textbf{255} (1979), 1--30.

\bibitem{Kalton}
N.~J. Kalton,
Twisted Hilbert spaces and unconditional structure,
\emph{J. Inst. Math. Jussieu} \textbf{2} (2003), 401--408.

\bibitem{Komorowski}
R.~Komorowski,
On constructing Banach spaces with no unconditional basis,
\emph{Proc. Amer. Math. Soc.} \textbf{120} (1994), 101--107.

\bibitem{MP}
B.~Maurey and G.~Pisier,
S\'eries de variables al\'eatoires vectorielles ind\'ependantes et propri\'et\'es
g\'eom\'etriques des espaces de Banach,
\emph{Studia Math.} \textbf{58} (1976), 45--90.

\bibitem{NTJ}
N.~J. Nielsen and N.~Tomczak-Jaegermann,
Banach lattices with property (H) and weak Hilbert spaces,
\emph{Illinois J. Math.} \textbf{36} (1992), 345--371.

\bibitem{Pisier}
G.~Pisier,
Weak Hilbert spaces,
\emph{Proc. London Math. Soc.} (3) \textbf{56} (1988), 547--579.

\bibitem{RK}
N.~de Rancourt and O.~Kurka,
The ergodicity of Orlicz sequence spaces,
\emph{J. Funct. Anal.} \textbf{290} (2026), no.~5, 111270.

\end{thebibliography}

\end{document}
